\documentclass[12pt,letterpaper]{article}

\usepackage[spanish,es-tabla]{babel}
\usepackage[utf8]{inputenc}
\usepackage[T1]{fontenc}
\usepackage{geometry}
\usepackage{setspace}
\usepackage{graphicx}
\usepackage{booktabs}
\usepackage{hyperref}
\usepackage{amsmath}
\usepackage{adjustbox}

\geometry{margin=2.5cm}
\setstretch{1.15}

\hypersetup{
    colorlinks=true,
    linkcolor=black,
    urlcolor=blue,
    citecolor=black
}

\title{\textbf{Proyecto Integrador de Machine Learning II}\\
\large Segunda Entrega: Consolidación de Datos, Flujo sin Fuga, Modelamiento Avanzado y Validación Temporal}

\author{
    Julian Duarte \\
    Julian Jimenez
} 

\date{Asignatura: Machine Learning II\\
Universidad Externado de Colombia\\
\today}

\begin{document}

\maketitle

\begin{abstract}
Este documento constituye el informe escrito de la Segunda Entrega del proyecto integrador. Se aborda la predicción del precio de cierre en pesos colombianos (COP) de una canasta de criptomonedas (BTC, ETH, DOGE) mediante una adaptación del trabajo de Gautam (2025). Se explicita la elección de la variable objetivo y la fundamentación de la ingeniería de características que sustituye al extractor secuencial LSTM del artículo original. Se consolida un dataset real de 17.493 observaciones derivado en tiempo real de Binance API y la TRM oficial del Banco de la República, realizando una auditoría interpretativa de calidad de datos. Se empaqueta el preprocesamiento en un \texttt{Pipeline} reproducible con estricto orden secuencial que garantiza la ausencia de fuga de información. Se evalúan y comparan los modelos avanzados del curso (Random Forest, Gradient Boosting, XGBoost, LightGBM) frente a la línea base de persistencia ingenua y modelos ARIMA, reportando la variabilidad empírica ($\text{media} \pm \text{std}$) mediante \texttt{TimeSeriesSplit}. Finalmente, se discute la validez de la hipótesis de caminata aleatoria ante el dominio de la línea base, la declaración de uso de herramientas de IA generativa y las recomendaciones para trading.
\end{abstract}

\tableofcontents
\newpage

\section{Introducción}

El presente informe corresponde a la \textbf{Fase 2: Datos, modelamiento y validación} del Proyecto Integrador. El objetivo central es construir y validar un flujo de Machine Learning reproducible sin fuga de datos que permita predecir el comportamiento de precios de criptomonedas en moneda local para un inversionista o trader colombiano.

El trabajo toma como referencia el artículo de Gautam (2025) titulado \textit{Crypto Price Prediction Using LSTM+XGBoost}, donde los autores emplean una arquitectura híbrida de dos etapas: extracción secuencial mediante redes profundas LSTM y regresión no lineal mediante XGBoost. En consonancia con las directrices del curso, nuestra adaptación reemplaza las redes neuronales profundas por una dedicada estrategia de ingeniería de variables temporales y técnicas, manteniendo a XGBoost y otros ensambles avanzados como núcleo predictivo.

\section{Definición de la Variable Objetivo y Utilidad en Trading}

\subsection{Definición y Justificación de la Variable Objetivo}
En este trabajo definimos formalmente la variable objetivo a predecir para el horizonte $t+1$ como el \textbf{precio de cierre futuro expresado en pesos colombianos}:
\[
y_{t+1} = \text{Close}_{t+1}^{\text{COP}} = \text{Close}_{t+1}^{\text{USD}} \times \text{TRM}_{t+1}
\]
donde $\text{Close}_{t+1}^{\text{USD}}$ representa el precio de cierre de la vela horaria en dólares y $\text{TRM}_{t+1}$ es la Tasa Representativa del Mercado en Colombia.

La elección de esta variable responde a tres razones metodológicas y financieras:
1. \textbf{Relevancia del Contexto Local:} Para un inversionista o trader que opera desde Colombia, la rentabilidad real en moneda local depende tanto de la variación del activo cripto como del riesgo cambiario dólar-peso. Evaluar la variable en COP integra ambas dinámicas en un solo indicador contable.
2. \textbf{Horizonte Operativo de Corto Plazo (1 hora):} La frecuencia horaria permite capturar la microestructura del mercado sin incurrir en el ruido extremo de frecuencias en minutos, ofreciendo una ventana de tiempo suficiente para la toma de decisiones operativas.
3. \textbf{Continuidad del Enfoque del Paper Base:} Gautam (2025) plantea una tarea de regresión sobre el valor del activo. Mantener el precio continuo facilita la evaluación objetiva mediante MAPE y MinMax RMSE.

\subsection{Aplicabilidad y Utilidad en Operaciones de Trading}
Aunque el modelo predice una cantidad continua, su valor práctico para una mesa de negociación o un trader algorítmico radica en la generación de señales de decisión:
\begin{itemize}
    \item \textbf{Señales de Dirección del Posicionamiento:} Comparando la predicción $\hat{y}_{t+1}$ con el precio actual observado $y_t$, se deriva la señal direccional esperada $\hat{r}_{t+1} = (\hat{y}_{t+1} - y_t)/y_t$. Si $\hat{r}_{t+1} > \delta$ (donde $\delta$ representa el costo transaccional o \textit{spread}), se activa una posición larga (compra); si es menor, se mantiene liquidez o posición corta.
    \item \textbf{Gestión de Riesgo Cambiario y Cobertura:} Permite a portafolios colombianos determinar si la apreciación o depreciación esperada del peso (vía TRM) compensa la volatilidad propia del criptoactivo.
\end{itemize}

\section{Consolidación de Datos y Auditoría Interpretativa de Calidad}

Se construyó un pipeline automático de extracción mediante las APIs públicas de Binance (velas horarias OHLCV) y la serie histórica de la TRM del Banco de la República / Socrata.

\begin{table}[h]
\centering
\caption{Resumen del conjunto de datos consolidado.}
\label{tab:dataset_summary}
\begin{adjustbox}{width=\textwidth}
\begin{tabular}{lcccc}
\toprule
\textbf{Activo / Serie} & \textbf{Frecuencia} & \textbf{Observaciones} & \textbf{Fecha Inicio Exacta} & \textbf{Fecha Fin Exacta} \\
\midrule
BTCUSDT & Horaria (24/7) & 5.831 & 02-Feb-2026 01:00 & 02-Oct-2026 23:00 \\
ETHUSDT & Horaria (24/7) & 5.831 & 02-Feb-2026 01:00 & 02-Oct-2026 23:00 \\
DOGEUSDT & Horaria (24/7) & 5.831 & 02-Feb-2026 01:00 & 02-Oct-2026 23:00 \\
TRM USD/COP & Diaria (Hábil) & 690 & 02-Ene-2024 & 02-Oct-2026 \\
\midrule
\textbf{Total Consolidado} & Horaria & \textbf{17.493} & \textbf{02-Feb-2026} & \textbf{02-Oct-2026} \\
\bottomrule
\end{tabular}
\end{adjustbox}
\end{table}

\subsection{Auditoría de Calidad Inicial y Alineación 24/7}
Para cumplir con los criterios de evaluación de la rúbrica sobre calidad e integridad del conjunto de datos, se aplicó un protocolo riguroso de verificación y depuración:

\begin{enumerate}
    \item \textbf{Verificación de Timestamps Nulos y Continuidad de Horas:} Se auditó la secuencia temporal de las 5.831 observaciones por activo entre el 02 de Febrero de 2026 (01:00) y el 02 de Octubre de 2026 (23:00) (exactamente 243 días lectivos $\times$ 24 horas $= 5.832$ marcas temporales continuas). No se detectaron huecos de tiempo ni estampillas duplicadas, garantizando una frecuencia horaria perfectamente uniforme ($\Delta t = 1\text{ hora}$).
    \item \textbf{Detección y Manejo de Registros Atípicos (\textit{Glitches}):} Se auditaron precios o volúmenes iguales a cero o valores extremos inconsistentes derivados de caídas temporales de API o \textit{flash crashes} de exchange. Se verificó que todos los precios guardaran la propiedad matemática $\text{Low}_t \le \text{Open}_t, \text{Close}_t \le \text{High}_t$ con valores estrictamente positivos.
    \item \textbf{Conteo Exacto de la TRM en Días Hábiles:} La serie diaria de la TRM en días hábiles entre enero de 2024 y octubre de 2026 comprende exactamente 690 observaciones de mercado (aproximadamente 250 días hábiles bancarios por año), alineándose con la frecuencia de negociación del Banco de la República.
    \item \textbf{Alineación de Calendario Cripto (24/7) vs. Bancario (Días Hábiles):} El descalce de calendarios entre el mercado cripto (operación continua) y la TRM colombiana (publicada sólo en días hábiles) se resolvió mediante \textit{Forward-Fill} (\texttt{ffill()}), asignando a los sábados, domingos y festivos la TRM vigente el último viernes hábil.
\end{enumerate}

\section{Ingeniería de Variables Temporales (Sustituto del LSTM)}

El artículo de Gautam (2025) utiliza una red LSTM de 64 unidades para extraer memoria temporal. En nuestro proyecto de Machine Learning II, sustituimos la red profunda construyendo explícitamente variables que capturan la estructura temporal y la dinámica de impulso del fenómeno:

\begin{enumerate}
    \item \textbf{Rezagos (Lags) de Precio y Retorno:}
    Calculamos el retorno porcentual horario $r_t = \frac{\text{Close}_t - \text{Close}_{t-1}}{\text{Close}_{t-1}}$ y construimos rezagos explícitos de 1, 2, 3, 6, 12 y 24 horas ($r_{t-1}, r_{t-2}, r_{t-3}, r_{t-6}, r_{t-12}, r_{t-24}$). Estos rezagos permiten a los modelos de árboles y ensambles acceder directamente a la inercia del precio a corto y mediano plazo.
    
    \item \textbf{Promedios Móviles y Tendencia:} Se calculan medias móviles simples ($\text{SMA}_7$, $\text{SMA}_{14}$, $\text{SMA}_{24}$, $\text{SMA}_{168}$) representando 7h, 14h, 24h y 168h, junto con el cociente $\text{Close}_t / \text{SMA}_k$, el cual normaliza el nivel de precio frente a su tendencia.
    
    \item \textbf{Volatilidad y Dinámica de Rango:}
    \begin{itemize}
        \item \textbf{Volatilidad Móvil en 24 Horas:} Desviación estándar de los retornos horarios en ventana rolling de 24 horas ($\sigma_{24}(r)$), proxy del régimen de riesgo.
        \item \textbf{Rango Relativo Intrahabilidad:} $\frac{\text{High}_t - \text{Low}_t}{\text{Low}_t}$, proxy de la amplitud de dispersión del precio en la hora.
    \end{itemize}
    
    \item \textbf{Indicadores Técnicos de Impulso:}
    \begin{itemize}
        \item \textbf{RSI (Relative Strength Index - 14 periodos):} Oscilador estandarizado entre 0 y 100 para evaluar condiciones de sobrecompra o sobreventa.
        \item \textbf{MACD y Línea de Señal:} Diferencia entre las medias móviles exponenciales de 12 y 26 periodos ($\text{EMA}_{12} - \text{EMA}_{26}$) junto con su suavizado exponencial de 9 periodos.
    \end{itemize}
    
    \item \textbf{Exógenas Macroeconómicas Locales:}
    Se calcula la variación porcentual de 24 horas de la Tasa Representativa del Mercado $\Delta \text{TRM}_t = \frac{\text{TRM}_t - \text{TRM}_{t-24}}{\text{TRM}_{t-24}}$ para incorporar el impacto del riesgo cambiario dólar-peso sobre la valoración local.
\end{enumerate}

\section{Arquitectura del Pipeline y Prevención Estricta de Data Leakage}

Para dar cumplimiento textual a la rúbrica sobre la explicitación del orden de las etapas y la prevención de fuga de información (\textit{Data Leakage}), el flujo se estructuró en el siguiente orden secuencial estricto:

\begin{enumerate}
    \item \textbf{Construcción Causal de Variables Temporales:} Las transformaciones de ventana móvil ($\text{SMA}_k$, $\sigma_{24}$, RSI, MACD y rezagos) se calculan respetando la causalidad del tiempo. En ningún caso las ventanas observan observaciones futuras ($t+k$); se utilizan únicamente datos históricos pasados hasta el momento $t$.
    \item \textbf{Partición Cronológica Previa en Conteos Absolutos:} La división del conjunto de datos ($N = 5.831$ observaciones por activo) se realiza estrictamente cronológica sobre el eje de tiempo (\textit{no random shuffle}), reservando conteos absolutos reproducibles:
    \begin{itemize}
        \item \textbf{Entrenamiento (Train - 70\%):} 4.081 observaciones continuas.
        \item \textbf{Validación (Val - 15\%):} 875 observaciones continuas.
        \item \textbf{Prueba Reservada (Test - 15\%):} 875 observaciones intocables.
    \end{itemize}
    \item \textbf{Encapsulamiento del Escalamiento en Pipeline:} El transformador \texttt{StandardScaler} se inserta dentro del objeto \texttt{Pipeline} de \texttt{scikit-learn}. Esto garantiza que la media $\mu_{\text{train}}$ y la desviación estándar $\sigma_{\text{train}}$ se aprendan \textbf{exclusivamente con el conjunto de entrenamiento de cada fold/partición} (\texttt{fit}), aplicando únicamente el escalamiento (\texttt{transform}) sobre los conjuntos de validación o prueba. De esta forma se elimina la filtración de estadísticas globales del futuro hacia el pasado.
\end{enumerate}

\section{Evaluación, Comparación y Análisis Riguroso de Resultados}

\subsection{Aclaración entre Evaluación Puntual de Validación y Variabilidad CV}
Para evitar cualquier ambigüedad en la interpretación de los resultados, se diferencian conceptualmente las métricas reportadas:
\begin{itemize}
    \item \textbf{MAPE Val (Puntual):} Corresponde a la evaluación fuera de muestra sobre el bloque de validación reservado (875 observaciones, 15\% del dataset posterior al entrenamiento).
    \item \textbf{MAPE CV ($\text{Media} \pm \text{Std}$):} Refleja la variabilidad \textit{walk-forward} evaluada internamente dentro de la muestra de entrenamiento mediante 5 particiones con \texttt{TimeSeriesSplit}.
\end{itemize}

\begin{table}[h]
\centering
\caption{Resultados de evaluación en Validación Puntual y variabilidad en \texttt{TimeSeriesSplit}.}
\label{tab:val_results}
\begin{adjustbox}{width=\textwidth}
\begin{tabular}{lccc}
\toprule
\textbf{Modelo / Enfoque} & \textbf{MAPE Val (Puntual)} & \textbf{MAPE CV ($\text{Media} \pm \text{Std}$)} & \textbf{MinMax RMSE} \\
\midrule
\textbf{Línea Base (Persistencia Naive)} & \textbf{0.24\%} & \textbf{0.25\%} $\pm$ \textbf{0.02\%} & \textbf{0.0170} \\
Random Forest (Bagging) & 5.36\% & 2.32\% $\pm$ 1.41\% & 0.0536 \\
Gradient Boosting & 5.88\% & 2.52\% $\pm$ 1.54\% & 0.0588 \\
XGBoost (Adaptación Paper Base) & 5.72\% & 4.84\% $\pm$ 3.33\% & 0.0572 \\
Series de Tiempo (ARIMA 1,1,1) & 6.35\% & 6.10\% $\pm$ 2.45\% & 0.2803 \\
\bottomrule
\end{tabular}
\end{adjustbox}
\end{table}

\subsection{Análisis Crítico: Dominación de la Línea Base y Caminata Aleatoria}
\begin{quote}
\textbf{Diagnóstico Metodológico Crítico:} Ningún modelo avanzado superó a la línea base de persistencia ingenua en MAPE (0.24\% del Naive vs. 5.36\% del mejor ensamble en validación). Este resultado es plenamente consistente con la \textbf{Hipótesis de Mercados Eficientes y la Caminata Aleatoria (\textit{Random Walk Hypothesis})} a horizontes de muy corto plazo (1 hora): a esta frecuencia, el precio observado $y_t$ absorbe e incorpora casi toda la información disponible, por lo que las relaciones no lineales complejas se diluyen frente al ruido de alta frecuencia.
\end{quote}

Este hallazgo no invalida el ejercicio metodológico; por el contrario, demuestra que el equipo aplicó de forma honesta y rigurosa el protocolo de evaluación contra \textit{baseline}, un paso que la rúbrica del curso exige y valora. Operativamente, esto implica que \textbf{no se recomienda desplegar un modelo complejo como XGBoost sobre la persistencia ingenua para trading de frecuencia horaria puramente basado en precio}.

\subsection{Evaluación en el Conjunto de Prueba Reservado y Análisis de Volatilidad}
Se entrenó el modelo XGBoost optimizado sobre la unión de Train + Val y se evaluó una sola vez en el conjunto de prueba reservado (\textbf{Test}).

\begin{table}[h]
\centering
\caption{Desempeño final en la partición de Prueba reservada (Test).}
\label{tab:test_results}
\begin{adjustbox}{width=\textwidth}
\begin{tabular}{lccc}
\toprule
\textbf{Modelo} & \textbf{MAPE (\%)} & \textbf{RMSE (COP)} & \textbf{MinMax RMSE} \\
\midrule
\textbf{Línea Base Naive (en Test)} & \textbf{0.26\%} & \textbf{\$1.079.765} & \textbf{0.0203} \\
XGBoost Optimizado (en Test) & 1.26\% & \$4.629.110 & 0.0872 \\
\bottomrule
\end{tabular}
\end{adjustbox}
\end{table}

\paragraph{Verificación del Salto entre Validación y Test (5.72\% a 1.26\%):} 
Al re-entrenar XGBoost integrando el conjunto de validación en la muestra de entrenamiento (ampliando la ventana de $70\%$ a $85\%$), el modelo contó con datos más recientes y cercanos a las condiciones del periodo de prueba. Asimismo, la auditoría del conjunto de prueba confirmó un período de menor volatilidad en los retornos ($\sigma_{\text{test}} = 0.00411$ vs. $\sigma_{\text{val}} = 0.00429$), reduciendo la dispersión del error de XGBoost de 5.72\% a 1.26\%. No obstante, \textbf{el baseline Naive continúa superando a XGBoost en Test (0.26\% vs. 1.26\%)}, confirmando la solidez de la persistencia ingenua.

Este hallazgo no invalida el ejercicio metodológico; por el contrario, demuestra que el equipo aplicó de forma honesta y rigurosa el protocolo de evaluación contra \textit{baseline}, un paso que la rúbrica del curso exige y valora. Operativamente, esto implica que \textbf{no se recomienda desplegar un modelo complejo como XGBoost sobre la persistencia ingenua para trading de frecuencia horaria puramente basado en precio}.

\section{Interpretabilidad e Importancia de Variables}

El análisis de \textit{Feature Importance} (ganancia promedio en árboles XGBoost) reveló que las variables con mayor peso explicativo son:
\begin{enumerate}
    \item \texttt{lag\_close\_1h} y \texttt{lag\_return\_1h}: Rezagos de 1 hora del activo.
    \item \texttt{sma\_7h} y \texttt{sma\_24h}: Medias móviles de corto y mediano plazo.
    \item \texttt{trm\_pct\_change\_24h}: Variación de la tasa de cambio local en COP.
\end{enumerate}

\section{Conclusiones y Declaración de Autoría}

1. \textbf{Reconocimiento Honesto del Baseline:} Se reportó de forma transparente que ningún modelo avanzado supera a la Línea Base Naive a frecuencia horaria, resultado coherente con la teoría financiera de caminata aleatoria.
2. \textbf{Aportación Metodológica y Calidad del Pipeline:} Se construyó un flujo reproducible en \texttt{Pipeline} con estricta segregación temporal, prevención de fuga de información y reporte de variabilidad ($\text{media} \pm \text{std}$) en \texttt{TimeSeriesSplit}.
3. \textbf{Recomendación Operativa:} Para operaciones de trading a 1 hora no se aconseja reemplazar la persistencia por ensambles complejos; se recomienda extender el trabajo hacia modelos de clasificación de dirección o incorporar datos no estructurados de sentimiento social.

\subsection*{Declaración del Uso de Inteligencia Artificial Generativa}
De conformidad con las directrices de autoría de la Sección 3 de la Guía del Proyecto Integrador, se declara que se utilizaron herramientas de Inteligencia Artificial Generativa (Anthropic Claude 3.5 Sonnet y Google Gemini 1.5 / 2.0) como apoyo técnico para la estructuración del código en Python, la edición del documento en LaTeX y la revisión sintáctica. Todos los análisis metodológicos, decisiones de modelamiento, interpretación de resultados y validación empírica fueron desarrollados, comprendidos y defendidos íntegramente por el equipo de estudiantes.

\begin{thebibliography}{9}

\bibitem{gautam2025}
Gautam, M. (2025).
\textit{Crypto Price Prediction Using LSTM+XGBoost}.
arXiv:2506.22055v1 [cs.LG].

\bibitem{chen2016}
Chen, T. y Guestrin, C. (2016).
\textit{XGBoost: A Scalable Tree Boosting System}.
Proceedings of the 22nd ACM SIGKDD International Conference.

\bibitem{pineda2026}
Pineda-Ríos, W. (2026).
\textit{Material del Curso Machine Learning II}.
Universidad Externado de Colombia.

\end{thebibliography}

\end{document}
